\hypertarget{local-services}{%
\chapter{7. Local services}\label{local-services}}

A workstation's most sensitive services do not listen on the network.
They listen locally, on sockets in the user's own session: the SSH and
GPG agents holding the keys, the display server, the audio server, the
session bus, the container daemon. These are the conveniences a desktop
is built from, and a workload running as the user can reach every one of
them.

For almost all of them, reaching is enough. Connection is authorization.
A workload that opens the SSH agent's socket can have the user's
challenges signed and authenticate as the user to every host that trusts
the key (\texttt{T3.7}); one that reaches the display server can read
the clipboard and watch the keystrokes (\texttt{T2.6}); the Docker
socket is, for practical purposes, root on the machine. There is no
second step where the service asks whether this particular client was
meant to have it. The socket is the authority, and under the default-uid
model every one of the user's sockets is open to anything running as the
user. The workload inherits the user's whole local session.

So the boundary cannot sit inside the conversation, asking whether a
given request to a service is allowed; by the time there is a
conversation, the authority is already handed over. The boundary has to
sit at the connection, and then, for the services where even an allowed
connection hands over too much, reach past it. What follows is that
progression: first making the connection a mediated thing at all, then
the graded set of answers for connections that cannot simply be granted.

\hypertarget{the-paths-then-the-connection}{%
\section{7.1 The paths, then the
connection}\label{the-paths-then-the-connection}}

The obvious way to confine local sockets is to keep doing what the
filesystem chapter did: build the workload a view with only the allowed
sockets in it, and let the absent ones be absent. It is the wrong tool
here, and the reason is worth stating, because it shapes what follows.

A socket is not data at rest. A socket path that appears in the
workload's world is not a fact it can read: it is an authority it can
exercise, out of band, the moment it sees it. So a view that contains
the allowed sockets still leaks: the workload reads its own grants off
what is present, and a present socket is a connectable one whether or
not the monitor was watching. The path model is brittle besides:
services that come and go through a session, sockets the host addresses
through rules the policy would have to mirror exactly, and names the
system resolves to a live endpoint though they were never files (an
abstract socket on one platform, a reserved DOS device name on another),
neither of which a path scheme can name at all. Each is a seam. Together
they say the filesystem is the wrong layer to hold this boundary.

So absence is applied to the paths wholesale. No host socket path is
present in the workload's world: not the forbidden ones, not the allowed
ones, none. There is nothing in the view to enumerate, probe, or connect
to behind the monitor's back, and where a name could reach an endpoint
without being a file the workload still has no way to address it. The
local surface is empty, the same way the network was.

What the workload has instead of paths is the monitor. A workload that
needs an allowed service asks for it by the name the policy gives it,
and the monitor, if the policy allows it, opens the real connection on
the host side and hands the connected channel back
(\texttt{interpose-as-transaction}). The application connects the way it
always has; the call becomes a request, and the request is answered or
refused. Enforcement moves off the filesystem and onto the connection,
where it can be checked the instant it is made and audited every time,
and where a service that restarted or moved since the workload began is
reconnected correctly, because the name is resolved fresh each time
rather than pinned to a path that may have gone stale.

\hypertarget{the-contract}{%
\section{7.2 The contract}\label{the-contract}}

The contract is default-deny, and the default holds by construction
rather than by rule: a service the policy does not name cannot be named
in a request, because the workload has no path to it and no handle on
it: only the set of names the policy gave it. There is nothing to add a
denylist to. The dangerous services (Docker, the system bus, the legacy
X11 server) are not refused; they are simply never named, and a service
that is never named is one the workload cannot ask for.

Asking by name rather than by path buys more than tidiness. The name is
the workload's only handle on the service: it has no path, no address,
nothing but the word the policy gave it, and so what stands behind that
word is the monitor's to choose. To the operator inspecting a kennel
before it runs, the choice is laid bare: the names the workload will ask
for, and what each one really reaches, side by side; the operator sees
the topology the workload never does. To the workload there is only the
name. Everything that follows is a different answer to one question:
when the workload asks for a service by name, what does the monitor
connect it to?

\hypertarget{the-raw-channel}{%
\section{7.3 The raw channel}\label{the-raw-channel}}

The simplest answer is the real service, connected raw. The monitor
opens the connection and hands over the channel, and from there it does
not listen: it granted the connection, it does not police the protocol
that runs across it (\texttt{mediate-use-not-reach}). A workload granted
the audio server can do anything PipeWire allows (play, and also
capture) and the monitor reads none of it.

This is the same edge the network boundary has, in the same place. That
a connection the policy forbids never opens is a safety property, and
the monitor enforces it completely: the forbidden connection is a thing
that, once it happens, cannot be called back, and the monitor sees that
it does not happen. What is said over a connection the policy allowed is
a different kind of thing (the use of the channel rather than the reach
to it), and a monitor watching one exchange at a time cannot bound it
{[}Schneider{]}. So a raw grant is a raw grant of whatever authority the
protocol confers: no more than was asked for, but no less.

Where the protocol's whole authority is what the operator meant to
grant, that is the right answer and the end of it. The audio server is
audio; granting it grants playing and capturing sound, which is what
audio access is. The answers that follow are for the services where the
protocol's whole authority is far more than anyone meant to hand across
in one piece.

\hypertarget{the-facade}{%
\section{7.4 The facade}\label{the-facade}}

The session bus is the clearest of those. A connection to D-Bus is a
connection to everything on it: the notification service, the secret
store, the network settings, the power controls: the desktop's whole
cast of services, reachable through a single socket. Handed the bus raw,
a workload has them all; connection is authorization, and here the
authorization is the session entire.

But the bus is not an opaque stream the way a signing agent's socket is.
It is structured (what travels on it is messages, calls to named
services and named methods), and structure is something a boundary can
hold. So where the bus cannot be handed over whole, it is not handed
over raw.

What the workload is given instead is a facade: something inside its own
kennel that wears the shape it expects (a socket that looks and answers
like the bus it meant to reach), so the application speaks to it
unaltered. The facade decides nothing. Its whole job is to present the
right face and carry what the workload sends across the kennel's IPC
surface to the other side, where the deciding happens: a separate
service kennel, or the host context, a place the workload has no handle
on. There the messages are parsed and policed, and only the calls the
policy named go on to the real bus. A workload granted the notification
service can post a notification; the message crosses to an enforcer it
cannot reach back into, and every call it was not granted stops there.

This is the shape of every brokered service, thrown into relief by the
danger of this one. The side facing the workload is only ever
presentation; the side that decides is always across the boundary,
beyond the workload's touch, which is what makes the deciding
trustworthy, because an enforcer the workload cannot reach is one it
cannot turn. It is also what keeps the trusted core small. Policing
D-Bus means understanding D-Bus, and understanding a protocol is a great
deal of code, exactly what a reference monitor is meant not to
accumulate. So it accumulates nowhere near the monitor: the protocol
logic lives in a confined service of its own, on the far side of the IPC
surface from the workload it guards against, and the monitor brokers the
connection and reads none of what crosses it
(\texttt{control-not-data-plane}). The workload faces a mask, and the
work goes on somewhere it cannot follow.

\hypertarget{the-private-instance}{%
\section{7.5 The private instance}\label{the-private-instance}}

Some services cannot be filtered into safety, because there is no clean
subset of the protocol to allow; and they cannot be handed over raw,
because raw is catastrophic; yet the workload has a real need the
service answers. The display is the standard case. A workload may
legitimately need to put a window on the screen; a connection to the
user's display server lets it read the clipboard, log every keystroke,
and capture the screen of every other window (\texttt{T2.6}). The
protocol is not a handful of discrete privileged calls to allowlist one
of; it is a continuous compositing relationship, with no message-level
line through it that leaves something both safe and useful.

So the design neither filters nor refuses. It substitutes. The workload
is given a display server of its own (a private, nested instance) and
composes into that. It draws its windows, receives the input meant for
it, and behaves exactly as it would against the real one, because as far
as it can tell it is talking to an ordinary display server. What it
cannot do is reach another window's input or output, because in its own
server there are no other windows. The user's real display is never
connected; the private instance stands on the far side of the boundary,
in its own confinement like every brokered service before it, so the
blast radius of a compromise is that instance and nothing past it.

The private instance is the answer; a raw connection to the real display
is the residual the answer exists to avoid, and it is worth naming,
because an operator can still ask for it. On Wayland, built with client
separation in mind, a raw connection to the real compositor still
carries the clipboard and the screen until a message-filtering facade
for it lands, so where a private instance is not in place that grant is
a real grant of that reach. On X11 it is past a residual. X11 has no
notion of one client kept from another; any client reads any other's
input and screen as a matter of protocol. There is nothing to filter and
nothing to contain, and a workload handed the user's real X server is
handed the whole session's eyes and ears. Granting it is equivalent to
turning the boundary off, and a workload that needs it is operating
outside the model.

\hypertarget{what-is-not-brokered}{%
\section{7.6 What is not brokered}\label{what-is-not-brokered}}

The last answer is no answer at all: the service is not brokered into
the kennel by any of the three routes. It is right for exactly one kind
of service: the one where there is no safe subset to filter, no useful
private instance to substitute, and raw access is the user's identity
itself.

The SSH and GPG agents are that kind. A socket to a signing agent is the
standing authority to authenticate as the user: to sign whatever
challenge is put to it, with no attestation of what is being signed or
to whom, no second check that the asker was meant to ask, and in some
protocols a path to the key material behind it. There is no safe subset
of ``sign this'' the framework could \emph{mediate}: a facade that
parsed the agent protocol and passed selected requests would be vouching
that a given signature is a legitimate one, which is a claim it has no
ground to make, and bounding which key the agent uses does not bound
what that signature then authenticates. And there is no useful private
instance, because the value of the agent is precisely the user's real
keys: a private agent holding none is just an absent one.

So the framework builds no signing facade: it will not stand a mediated
service in front of an agent and vouch for what crosses it. This is a
bound on what the framework constructs, not on what the operator may
reach. The raw channel of §7.3 remains available for an agent socket as
for any other, and an operator who writes that grant has connected the
workload to the agent with the framework parsing nothing and vouching
for nothing; it is a warned footgun the operator owns, loud and tagged,
not a sanctioned path. What the framework refuses is to \emph{be} the
thing that mediates the signing, never to let the operator express a raw
connect. The legitimate need behind the SSH agent, a workload that must
reach a host over SSH, is real, and it is met a better way than the raw
socket: a credential bound to a single destination, with the real key
kept on the host side and out of the kennel entirely. That is a boundary
of its own shape, standing on the network rather than the socket, and it
is the subject of the next chapter. The GPG case has no such
re-originated form, for the reason the next chapter draws out: it is
attestation, and attestation does not belong in confinement at all.

The two limbs meet here the way they met at the network, and then
interposition goes further than it did there. Absence takes the paths,
so there is no service to reach except by asking. Interposition answers
the asking, and because the workload's only handle is a name, the
monitor is free in what it puts behind that name: the real service raw,
a facade that filters it, a private instance that stands in for it, or
nothing at all. Connection was the authorization, so the design made the
connection the transaction, and then made the transaction answerable in
whatever measure the service's danger demands. Each rung leaves a
residual, named where it falls (the protocol's full authority where it
is granted raw, the correctness of the enforcer behind a facade where it
is filtered, the blast radius of the private instance where one stands
in), and the service whose only honest answer was refusal is handed to
the chapter that can give it one.
